\chapter{Guía de desarrollo}
\label{cap:desarrollo}

\section{Estructura de repositorios}

El ecosistema Atlas se organiza en dos familias de repositorios:

\begin{itemize}
\item \textbf{Servicios de negocio} (prefijo \texttt{atlas-}): API, CMS, scraping y capacidades IA.
\item \textbf{Librerías de infraestructura} (prefijo \texttt{ferimer-}): colas, endpoints y acceso a base de datos.
\end{itemize}

Convenciones recomendadas:

\begin{itemize}
\item Nombres de repositorio en minúsculas y separados por guiones.
\item Módulos desacoplados por responsabilidad única.
\item Contratos JSON y documentación técnica versionados en el mismo ciclo que el código.
\item Dependencias compartidas centralizadas en librerías comunes para evitar duplicidad de lógica.
\end{itemize}

\section{Añadir un nuevo endpoint a la API}

Para incorporar un endpoint nuevo en la API Gateway se recomienda el siguiente flujo:

\begin{enumerate}
\item Definir el contrato del endpoint (método HTTP, ruta, parámetros y códigos de respuesta).
\item Asociar la operación a un tag de ontología con formato \texttt{dominio:categoría:acción[:subtipo]}.
\item Añadir/actualizar los esquemas \texttt{request} y \texttt{response} en \texttt{atlas-messaging-ontology}.
\item Configurar el endpoint declarativo en la API para mapearlo a su \textit{routing key} AMQP.
\item Activar validación AJV para entrada/salida y verificar errores de contrato.
\item Probar el flujo completo en modo RPC verificando \texttt{correlation\_id} y \texttt{reply\_to}.
\end{enumerate}

Este procedimiento asegura consistencia entre interfaz HTTP, mensajería interna y validación de datos.

\section{Añadir un nuevo scraper}

Los scrapers se extienden a partir de un contrato base común para facilitar reutilización y trazabilidad.

Pasos recomendados:

\begin{enumerate}
\item Implementar un nuevo adaptador extendiendo \texttt{BaseScrapper}.
\item Reutilizar \texttt{ICSScrapper} o \texttt{FBScrapper} como referencia cuando la fuente sea similar.
\item Normalizar campos obligatorios (título, fecha, ubicación, fuente, URL, idioma).
\item Gestionar deduplicación y control de fechas para evitar eventos repetidos.
\item Publicar resultado en formato uniforme para su validación y consumo posterior.
\item Registrar logs de extracción para diagnóstico de cambios en la fuente origen.
\end{enumerate}

Cada nuevo scraper debe poder activarse o desactivarse sin impactar al resto del pipeline.

\section{Añadir un nuevo servicio RabbitMQ}

La incorporación de un nuevo servicio AMQP requiere sincronizar implementación, contratos y enrutado.

\begin{enumerate}
\item Definir el nuevo dominio funcional y su tag de ontología.
\item Crear esquemas \texttt{request/response} en \texttt{atlas-messaging-ontology}.
\item Implementar un \textit{receiver} usando \texttt{AMQP\_RPC\_Receiver} o \texttt{AMQP\_FireAndForget\_Receiver}.
\item Declarar exchange, cola y binding con \textit{routing key} correspondiente.
\item Validar payloads entrantes y respuestas contra los esquemas oficiales.
\item Probar interoperabilidad desde \textit{emitters} de API u otros microservicios.
\end{enumerate}

El servicio debe mantenerse stateless y sin dependencias implícitas de orden temporal con otros consumidores.

\section{Convenciones de código}

Para mantener coherencia entre repositorios se establecen las siguientes convenciones:

\begin{itemize}
\item Módulos JavaScript en formato ESM.
\item Formateo automático con Prettier en CI y en desarrollo local.
\item Reglas de calidad y validación previas a commit mediante \textit{pre-commit hooks}.
\item Mensajes de commit orientados a cambio funcional y versionado semántico para librerías compartidas.
\item Separación clara entre configuración por defecto y secretos de entorno.
\end{itemize}

La aplicación consistente de estas pautas reduce deuda técnica y facilita mantenimiento transversal.

\section{Tests}

La estrategia de pruebas combina validación unitaria, integración y carga.

\begin{itemize}
\item \textbf{Unitarias}: lógica de transformación, utilidades y adaptadores aislados.
\item \textbf{Integración}: contratos AMQP, validación de esquemas y conectividad con servicios de infraestructura.
\item \textbf{End-to-end}: recorrido completo API $\rightarrow$ RabbitMQ $\rightarrow$ microservicio $\rightarrow$ respuesta.
\item \textbf{Carga}: pruebas de rendimiento con \texttt{k6} para medir latencias, throughput y comportamiento bajo estrés.
\end{itemize}

En el ecosistema Node.js pueden coexistir el \textit{Node.js test runner} y Jest según la madurez o naturaleza de cada repositorio, siempre con cobertura mínima definida por equipo y ejecución automatizada en CI.
